% ===== Appendix A: geometric phase, technical sketch =====
\appendix
\section{Geometric Phase: A Technical Sketch}
\label{p10:app:phase}

\noindent
This appendix records, in compressed form, the formal apparatus to which the main text's geometric language points. It is a sketch, not a proof, and it is deliberately so: the conjecture of \xref{p10:sec:invariant}, that holonomy is the gauge-invariant core conserved across the symmetry-broken phases, is offered as unsettled, and this appendix supplies its formal vocabulary without pretending to settle it. The Value-Foam programme is the place where the carrier of the phase must be given its independent quantitative content; here we record only the structural skeleton, with three honest cautions appended.

\subsection*{A.1\quad The bundle}
Let $B$ be the \emph{base}: the space of relational situations or contents, the configurations a relation may occupy. Over each point of $B$ stands a \emph{fibre} $F$: the phase of value or recognition (following the main text's anchoring, \xref{p10:sec:identity}, the sublimation of jouissance): the carrier on which the good circle's surplus is borne. A relation's history is a path in $B$; the parallel transport of its state in $F$ along that path is governed by a connection $\omega$. The phase accrued on returning to a starting point in $B$, the failure of the transported fibre-element to coincide with its origin, is the \emph{holonomy} of the loop,
\[
\holo(\gamma) \;=\; \mathcal{P}\exp\!\oint_\gamma \omega .
\]

\subsection*{A.2\quad The good and the vicious circle, formally}
A \emph{vicious} circle is a loop of vanishing holonomy: the relation returns in $B$ to where it began and the fibre returns unchanged (pure repetition, the bad infinite, the circle that goes round without rising, at zero phase); or a loop of \emph{negative} holonomy (net extraction, the fibre depleted on return). A \emph{good} circle is a loop of \emph{positive} holonomy: a return to the root in $B$ that nonetheless accrues a non-trivial, irreducible phase in $F$ (the spiral; sublimation; the bringing-back of a remainder). The criterion of \pinkref{Paper~IX} is the sign of $\holo(\gamma)$.

\subsection*{A.3\quad Variable basis as variable connection}
The main text's claim (\xref{p10:sec:measurement}) that the manner of attending is the choice of a measurement basis has its formal image here: a change in the manner of attending is a change of gauge, altering $\omega$, hence altering the holonomy a given loop in $B$ will accrue. Two relations may traverse the same situations (the same $\gamma$ in $B$) and accrue opposite phases, because they attend (connect) differently. This is why no fact about the base alone (the situations, the ``what happened'') fixes the goodness of a circle; the connection, the manner, is constitutive.

\subsection*{A.4\quad Personal identity as parallel transport}
The subject's continuity through the reflexive cycle (\xref{p10:sec:identity}) is parallel transport in this bundle: the same subject carried around a loop, returning to its root in $B$ bearing a phase in $F$. Alienation is negative-phase drift (the subject returns diminished, not recognising itself); self-realisation is positive-phase sublimation (the subject returns raised, more itself). ``The creator is the beneficiary'' is sameness in the holonomic sense, identity that has accrued a phase.

\subsection*{A.5\quad Three cautions}
The sketch must not be mistaken for more than it is.

\emph{First, the carrier.} A geometric phase requires a genuine bearer with the right structure (an $S^1$ or group-valued quantity, additive, gauge-invariant in the holonomy), failing which the apparatus is mere figure. The main text anchors the carrier in the sublimation of jouissance (\xref{p10:sec:identity}); whether that quantity in fact carries the required structure is exactly what the Value-Foam programme must establish independently, and is not settled here.

\emph{Second, sublimation and sign.} That ``sublimation $=$ positive phase'' and ``depletion $=$ negative phase'' must be \emph{argued}, not stipulated as a definitional identity, on pain of circularity. The independent route is to define a reproduction-surplus, net regeneration of the relation's generative conditions per cycle, minus consumption, and then show that its sign tracks the holonomy's; absent that, the identification is a picture, not a result.

\emph{Third, the cross-phase conjecture.} The proposal that $\holo$ is \emph{the} invariant conserved across all the symmetry-broken disciplinary phases (\xref{p10:sec:invariant}) is the most speculative claim in the series, and the appendix does nothing to discharge it. It records the conjecture's formal shape (an invariant belonging to no chart, conserved under all transition functions) and leaves it, as the main text insisted it must be left, offered and unsettled, to be walked rather than declared.
